% Shared preamble for the supporting-information documents (S1-S5).
\documentclass[10pt,letterpaper]{article}
\usepackage[margin=1in]{geometry}
\usepackage{amsmath,amssymb}
\usepackage{booktabs,array,longtable}
\usepackage[table]{xcolor}
\usepackage{hyperref}
\usepackage{microtype}
\usepackage{listings}
\usepackage{graphicx}
\usepackage[labelfont=bf,labelsep=period,justification=raggedright,singlelinecheck=off]{caption}
\definecolor{linkblue}{HTML}{1C5CAB}
\definecolor{codebg}{HTML}{F4F4F1}
\hypersetup{colorlinks=true, linkcolor=linkblue, citecolor=linkblue, urlcolor=linkblue}
\lstset{basicstyle=\ttfamily\footnotesize, backgroundcolor=\color{codebg},
  frame=none, breaklines=true, columns=fullflexible, keepspaces=true,
  xleftmargin=4pt, xrightmargin=4pt, aboveskip=6pt, belowskip=6pt}
\setlength{\parindent}{0pt}
\setlength{\parskip}{5pt}
\raggedright
\renewcommand{\tablename}{Table}
\newcommand{\sititle}[2]{{\Large\bfseries #1}\par\vspace{3pt}{\large #2}\par\vspace{2pt}
  {\small\color{black!60}Supporting information for ``Environmental fluctuation limits gene regulation unless it can be tracked''}\par\vspace{10pt}}

\begin{document}
\sititle{S1 Text.}{ODD description of EvoSim}

This description follows the ODD protocol in its 2020 revision (Grimm et al.,
J Artif Soc Soc Simul 23(2):7). It documents the version of EvoSim used for the
experiments in the main text, and every parameter value is given in
Table~\ref{tab:params}. Where a mechanism is switched off in these experiments it
is described briefly and marked as unused.

\section*{1. Purpose and patterns}
EvoSim is an individual-based model of unicellular organisms whose genomes are
gene regulatory networks of evolvable size and connectivity. In this study it is
used to ask how environmental fluctuation, and the information available to track
it, affect the number of regulatory connections and their strength. The model is built so that network size and connectivity evolve under mutation, selection and drift. Each phenotypic trait has a cost as well as a benefit, so traits settle at interior values, and the resource ecology is built partly by the organisms themselves through the byproducts of consumption.

\section*{2. Entities, state variables and scales}

\paragraph*{Grid cells.} The environment is a square lattice of $G\times G$ cells
($G=128$ in this study) with reflecting boundaries. Each cell holds the concentration of
six resource channels, $R_{r}(x,y)$ for $r=1,\dots,6$, capped at 6.0, and a count of
the organisms occupying it (at most 8).

\paragraph*{Organisms.} Each organism has a position, an energy store (at most
4.0), an age in time steps, a lineage depth (number of ancestral divisions), a
deleterious-load counter, and a genome. The genome is a regulatory network with
$n_g$ nodes, $13 \le n_g \le 32$, consisting of a weight matrix
$W\in\mathbb{R}^{n_g\times n_g}$ and a bias vector $b\in\mathbb{R}^{n_g}$. Nodes
1--12 are the phenotype outputs, node 13 is the sensor for the abiotic condition,
and nodes $14,\dots,n_g$ are hidden nodes with no fixed meaning. Each organism
also carries its phenotype, computed once at birth and fixed for life in these experiments. It consists of twelve traits in $[0,1]$, namely six resource-uptake shares (normalised to sum to one), body size,
motility, division threshold, DNA-repair fidelity, stress tolerance and leakage.

\paragraph*{Global state.} The abiotic condition, which cycles through discrete
states (Section 7.3); the seasonal phase; and a famine counter.

\paragraph*{Scales.} Time advances in discrete steps. In the experiments of this study one generation took about 24 time steps in fluctuating and 31 in frozen
populations (measured from the slope of mean lineage depth over the second half of
each run), so a 150{,}000-step run covers roughly 5{,}000 to 6{,}000 generations.
The maximum age is 60 steps.

\section*{3. Process overview and scheduling}
Each time step executes the following steps in order.
\begin{enumerate}
\item \emph{Resource renewal.} Every channel in every cell decays, and external
  inflow is added (Section 7.1).
\item \emph{Shock check.} With a small probability a mass-mortality or famine event
  occurs (Section 7.9).
\item \emph{Reversible development} (Section 7.4), which was not used in this study.
\item \emph{Organism update.} Organisms are visited once each, in storage order.
  For each organism the step comprises uptake and conversion of resources, with byproduct release into the
  cell (Section 7.5); maintenance cost (7.6); energy update; ageing unless
  dormant (7.7); movement (7.8); death if energy is exhausted or the maximum age is
  exceeded; otherwise, if the energy store exceeds the division threshold and the
  organism is not dormant, it is queued to divide.
\item \emph{Births} (Section 7.10). Queued parents are processed in queue order.
  Placement is serial; mutation and development of the offspring are computed in
  parallel with a random seed drawn serially for each offspring, so results do not
  depend on the number of threads.
\item \emph{Diffusion} of all channels (Section 7.2).
\end{enumerate}
Because organisms are updated sequentially and draw on shared cell resources, an
organism visited earlier in a step can deplete a cell before a later one reaches
it. Storage order reflects the reuse of memory slots and has no biological meaning.

\section*{4. Design concepts}
\paragraph*{Basic principles.} Every trait that brings a benefit also carries a
cost, so that selection does not drive traits to a boundary. Uptake is an allocation of a fixed budget across the six channels, so the affinities are not independent. Consumption of channel $r$ releases a byproduct into channel $r+1$,
so all channels except the first exist only through the activity of the
organisms.
\paragraph*{Emergence.} Network size, connectivity and connection strength,
population size, resource distribution and the composition of the community
emerge from the individual processes.
\paragraph*{Adaptation.} Organisms do not adapt within their lifetime. The
phenotype is fixed at birth and is a function of the genome, the local resource
concentrations and the value of the abiotic cue at the place and time of birth.
\paragraph*{Objectives, learning and prediction.} None; there is no explicit
fitness function. Reproductive success follows from energy balance.
\paragraph*{Sensing.} At birth, an organism senses the concentrations of the six
channels in its cell and the value of the abiotic cue.
\paragraph*{Interaction.} Organisms interact through shared resources, through
byproducts that others can consume, through crowding costs among organisms in the
same cell, and through the occupancy limit that can block a division.
\paragraph*{Stochasticity.} Mutation, movement, the placement of offspring,
shocks and the founding genomes are random. All randomness derives from the
simulation seed, and runs are exactly reproducible for a given seed.
\paragraph*{Collectives.} None are modelled explicitly. Lineages are tracked as
clades for bookkeeping only; they do not affect the dynamics.
\paragraph*{Observation.} Population means of the phenotype, of network structure
(number of nodes, number of non-zero weights, mean absolute weight) and of
ecological quantities are recorded at every time step. The measures used in the
main text are defined in Section 7.12.

\section*{5. Initialisation}
All cells start with $R_1=1.0$ and the other channels empty. $\max(200, G^2/80)$
founders (204 at $G=128$) are placed in uniformly random cells with energy 0.8,
age 0 and load 0. Each founder has $n_g=16$ nodes; each weight is non-zero with
probability 0.25 and then drawn from $\mathcal{N}(0,0.6^2)$, and biases are drawn
from $\mathcal{N}(0,0.2^2)$. Founders are developed in the conditions of their
cell before the first step.

\section*{6. Input data}
None. The model uses no external data.

\section*{7. Submodels}

\subsection*{7.1 Resource renewal}
Each step, every channel in every cell is multiplied by $1-0.006$. External inflow
per cell is
\begin{equation*}
I(t) = 0.25\,\Big(1 + 0.6\sin\frac{2\pi t}{15{,}000}\Big)\times f(t),
\end{equation*}
where $f(t)=0.65$ during a famine and 1 otherwise, and $I(t)$ is set to zero if
negative. In frozen environments (one condition state) all inflow enters channel
1. In fluctuating environments 60\% enters channel 1 and 40\% enters the channel
with the highest yield in the current state (Section 7.3). Concentrations are
capped at 6.0 after each addition.

\subsection*{7.2 Diffusion}
Each channel diffuses by a five-point Laplacian with rate 0.12 per step and
reflecting (no-flux) boundaries; concentrations are capped at 6.0.

\subsection*{7.3 Abiotic condition, state and cue}
The condition has $n_s$ states ($n_s=6$ in fluctuating arms, 1 in frozen arms) and
a period of $P=2{,}000$ steps. With phase $\theta(t) = (t \bmod P)/P$, the current
state is $s = 1 + \lfloor n_s\theta\rfloor$. The state rotates the yields, so that the yield of channel $r$ in state $s$ is $y_{1 + ((r-1+s-1) \bmod 6)}$, with
$y = (1.00, 0.60, 0.45, 0.35, 0.28, 0.22)$, so a different channel pays best in each state. The cue delivered to node 13 is defined as follows.
\begin{itemize}
\item In the informative arm $c = 2\theta - 1$, which determines the state exactly.
\item In the yoked arm $c = 2\phi - 1$, with $\phi$ the fractional part of
  $(t \bmod P)\cdot 6.1803399/P$. The multiplier $6.1803399 = 10/\varphi$, with
  $\varphi$ the golden ratio, is irrational, so the cue covers the same range
  without locking to the cycle and carries 0.0004 bits of information about the
  state.
\item In the no-cue arm $c=0$.
\end{itemize}
With $n_s = 1$ the state is always 1 and the cue is zero, so the three frozen arms
are identical.

\subsection*{7.4 Development}
At birth the expression vector is initialised to $e_k = 0.1$ and updated
synchronously eight times,
\begin{equation*}
e_k \leftarrow \tanh\Big(b_k + \sum_{m=1}^{n_g} W_{mk}\,e_m + u_k\Big),
\end{equation*}
where $u_k = R_k$ at the organism's cell for $k\le 6$, $u_{13}=c$, and $u_k=0$
otherwise. The phenotype is $\phi_k = (e_k+1)/2$ for $k=1,\dots,12$, and the
uptake shares $\phi_1,\dots,\phi_6$ are rescaled to sum to one. Development does
not reach a fixed point in general; the phenotype is the state after the eighth
update, which is deterministic and heritable. Reversible development, in which living organisms are re-developed every few steps, exists in the model but was not used in this study.

\subsection*{7.5 Uptake and conversion}
Let $\ell = 0.45 + 0.55\,\phi_{12}$ be the effective leakage and
$m = (0.5+\phi_7)(1+1.2\,\ell)$ the rate multiplier (multiplied by 0.35 when
dormant). For each channel $r$ with share $a_r=\phi_r>0$ and local concentration
$R_r$, the amount taken is $g_r=\min(0.6\,R_r\,a_r\,m,\;R_r)$. It is removed from
the cell, the organism gains $g_r\,y'_r\,(1-\ell)$ energy, where $y'_r$ is the
current yield of channel $r$, and $g_r\,\ell$ is added to channel $r+1$ in the same
cell (the byproduct of channel 6 leaves the system). The summed concentration of
channels 2--6 in the cell is recorded as the local toxicity $T$.

\subsection*{7.6 Maintenance}
The energy cost per step is
\begin{multline*}
0.08 + 0.22\,\phi_7^2 + 0.035\,\phi_8 + 0.04\,\phi_{11}^2 + 0.02\,\phi_{10}
+ 0.0015\,(n_g-13) + \kappa_e\, n_{nz}\\
+ 0.05\,(\text{other organisms in the cell}) + 0.015\,L
+ 0.008\,T\,(1-\phi_{11}),
\end{multline*}
where $n_{nz}$ is the number of non-zero weights, $\kappa_e = 5\times10^{-5}$ in
these experiments, and $L$ is the load counter. The cost is multiplied by 0.05
when dormant. Energy is updated as gains minus cost, capped at 4.0.

\subsection*{7.7 Dormancy and ageing}
An organism whose energy at the start of the step is below 0.15 is dormant for that step. A dormant organism has reduced uptake (Section 7.5) and a reduced maintenance cost (Section 7.6), does not age and cannot divide. Otherwise its age increases
by one.

\subsection*{7.8 Movement}
With probability $\phi_8$ the organism moves to a cell chosen uniformly from its
Moore neighbourhood including its own cell, clamped at the boundary.

\subsection*{7.9 Shocks}
At each step, with probability $1/25{,}000$, one of three events occurs, each with
probability 1/3: a cull, in which every organism dies with probability 0.92; a
heat shock, in which each organism dies with probability $0.92\,(1-\phi_{11})$;
or a famine, which multiplies inflow by 0.65 for the next 400 steps. These events
occur in all arms.

\subsection*{7.10 Division and inheritance}
An organism that is not dormant and whose energy is at least $2.0\,(0.5+\phi_9)$
divides. The offspring is placed in a random cell of the parent's Moore
neighbourhood (including the parent's cell); if that cell already holds 8
organisms the division fails. Parent and offspring each receive half the parent's
energy. The offspring has age 0, lineage depth one greater than the parent, and
load equal to the parent's load plus the change from mutation (Section 7.11),
floored at zero. Its genome is a mutated copy of the parent's, and it is
developed in the conditions of its own cell.

\subsection*{7.11 Mutation}
The per-site rate is $\mu = \max\big(10^{-3}(1 - 0.85\,\phi_{10}^{\text{parent}}),\,10^{-6}\big)$,
so DNA repair can reduce it by up to 85\%. At each birth the offspring genome
receives a Poisson number of point mutations with mean $\mu n_g^2$, each adding
$\mathcal{N}(0,0.25^2)$ to a randomly chosen weight, and with probability $\mu n_g$
one randomly chosen bias receives $\mathcal{N}(0,0.25^2)$ as well. With
probability 0.2 (in these experiments; 0 by default) one randomly chosen weight is
set to exactly zero. We call this connection silencing. Without it, additive
mutation would drive every weight away from zero and density towards one.

Structural mutations act on whole nodes. With probability 0.002 all outgoing
weights of a random node are set to zero (gene silencing). With probability 0.002
the incoming and outgoing weights and the bias of a random node are copied onto
another random node, and $\mathcal{N}(0,0.05^2)$ noise is added to every copied
weight, including weights that were zero in the source (gene duplication). With
probability 0.003 a hidden node is added, if $n_g<32$, as a copy of a random
existing node made in the same way as a duplication. If no node is added, a random
hidden node is removed with probability 0.003, provided $n_g>13$.
In addition, each point mutation, bias mutation and connection silencing is
classified for the load counter as deleterious (probability 0.70, load $+1$),
beneficial (0.05, load $-1$) or neutral (otherwise), and each gene silencing adds
1 to the load. Load carries a maintenance cost (Section 7.6) and can fall as well
as rise.

\subsection*{7.12 Measures used in the main text}
Two quantities are computed for each organism at birth, the number $n_{nz}$ of weights with $|w|>10^{-4}$ and the strength density $\sum|w|/n_g^2$ over all $n_g^2$ weights, zeros included. At each time step the population means of $n_{nz}$, of strength density and of $n_g$ are recorded. Count density is the mean $n_{nz}$ divided by the squared mean $n_g$, and total strength $\sum|w|$ is the mean strength density times the squared mean $n_g$. Each run is summarised by the mean over its final 20\% of time
steps, with ratios computed per step before averaging.

\section*{Parameters}
\begin{longtable}{@{}p{0.34\linewidth}p{0.22\linewidth}p{0.34\linewidth}@{}}
\caption{EvoSim parameters. ``Default'' is the value in the code; ``This study''
lists values that differ from the default.}\label{tab:params}\\
\toprule
Parameter & Default & This study\\
\midrule
\endfirsthead
\toprule
Parameter & Default & This study\\
\midrule
\endhead
Grid size $G$ & 256 & 128\\
Maximum organisms per cell & 8 & \\
Founders & $\max(200, G^2/80)$ & 204\\
Base inflow & 0.25 & \\
Season period, amplitude & 15{,}000; 0.6 & \\
Decay per step & 0.006 & \\
Diffusion rate & 0.12 & \\
Resource cap & 6.0 & \\
Yields $y_1,\dots,y_6$ & 1.00, 0.60, 0.45, 0.35, 0.28, 0.22 & \\
Condition states $n_s$ & 1 & 6 (fluctuating), 1 (frozen)\\
Condition period $P$ & 2{,}000 & \\
Inflow share following the favoured channel & 0 & 0.4 (fluctuating)\\
Cue gain & 1 & 0 (no-cue arms)\\
Yoked cue & off & on (yoked arms)\\
Uptake rate & 0.60 & \\
Leakage floor; rate bonus & 0.45; 1.2 & \\
Maintenance per step (base, size$^2$, motility, stress$^2$, repair) & 0.08, 0.22, 0.035, 0.04, 0.02 & \\
Maintenance per hidden node & 0.0015 & \\
Maintenance per non-zero weight $\kappa_e$ & 0 & $5\times10^{-5}$\\
Maintenance per step (crowding, load, toxicity) & 0.05, 0.015, 0.008 & \\
Division energy & 2.0 & \\
Initial energy; maximum energy & 0.8; 4.0 & \\
Dormancy threshold; upkeep; uptake & 0.15; 0.05; 0.35 & \\
Maximum age & 60 & \\
Development iterations & 8 & \\
Per-site mutation rate; $\sigma$ & $10^{-3}$; 0.25 & \\
Maximum repair suppression & 0.85 & \\
Gene silencing; duplication & 0.002; 0.002 & \\
Node addition; deletion & 0.003; 0.003 & \\
Connection silencing per birth & 0 & 0.2\\
Deleterious; beneficial mutation share & 0.70; 0.05 & \\
Initial nodes & 16 & \\
Shock probability per step; kill fraction & $1/25{,}000$; 0.92 & \\
Famine length; inflow factor & 400; 0.65 & \\
Run length & & 150{,}000 steps\\
Seeds & & 1--30 per arm\\
\bottomrule
\end{longtable}

\end{document}
